\documentclass{article}
\usepackage{kotex}
\usepackage{setspace}  % 줄간격
\setstretch{1.15}
\usepackage{enumitem}  % 목록 들여쓰기 
\setlist[itemize]{leftmargin=*, align=left}
\setlist[description]{
  font=\normalfont,
  labelwidth=4cm,
  leftmargin=4cm,
  labelsep=0cm,
  align=left
}
\usepackage{titlesec}  % 섹션 스타일 조정
\titleformat{\section}
  {\normalfont\Large\bfseries}
  {}
  {0pt}
  {}
  [\vspace{0.3em}\titlerule]

\titlespacing*{\section}
  {0pt}
  {1.5em}
  {0.8em}
\newcommand{\descriptionrule}{%
  \item[]\vspace{-0.6em}
  \noindent\hspace*{-\leftmargin}%
  \rule{0.6\linewidth}{0.4pt}
  \vspace{-0.4em}
}              
  \usepackage[a4paper, margin=2.5cm]{geometry}
\usepackage{fancyhdr}                 
\pagestyle{fancy}
\fancyhf{}                             
\fancyfoot[L]{\nouppercase{통사론 (2026학년도 2학기)}} 
\fancyfoot[R]{\thepage}                
\renewcommand{\headrulewidth}{0pt}    
\renewcommand{\footrulewidth}{0.4pt}   

\begin{document}

\title{통사론 (2026학년도 2학기)}
\author{\textit{written by using} \LaTeX}
\date{\today}
\maketitle
\thispagestyle{fancy}
\setcounter{secnumdepth}{0}
\bigskip

\section{0. 강좌정보}
\medskip
교과목번호: 108.315(001)\\
개설학과: 서울대학교 인문대학 언어학과\\
담당교수: 강초롱\\
수업시간: 월/수 11:00--12:15\\
수업장소: 14-602

\section{1. 수업목표}
\medskip
이 수업은 생성문법을 기반한 학부 통사론 교과목으로, 학생들은 생성문법의 기본 가정과 이론들을 배우게 됩니다. 영어를 기본으로 다양한 인간 언어에서 발견되는 통사 구문을 분석하는 방법을 배우고, 통사론에서 사용하는 용어와 이론들에 익숙해지며 나아가 통사적 논증방법을 배웁니다.

\section{2. 교재 및 참고문헌}
\medskip
\subsection{교재}
『Syntax: A Generative Introduction(4th Edition.)』(Andrew Carnie, 2021.)

\newpage
\section{3. 주차별 강의계획\protect\footnotemark}
\footnotetext{수강신청 사이트 `강의계획서' 항목 참조}
\medskip
\begin{description}
    \item[1주차(9/2, 9/7)] Generative Grammar (Chapter 1), Parts of speech (Chapter. 2)
    \item[2주차(9/9, 9/14)] Constituency, Trees, and Rules (Chapter. 3)
    \item[3주차(9/16, 9/21)] Structural Relations (Chapter. 4) 
    \item[4주차(9/23, 9/28)] Binding Theory (Chapter 5)
    \item[5주차(9/30)] X-bar Theory (Chapter 6)
    \item[6주차(10/7, 10/12)] Extending X-bar Theory to Functional Categories (Chapter 7)
    \item[7주차(10/14, 10/19)] Constraining X-bar: Theta Theory (Chapter 8)
    \item[8주차] Review(10/21), Midterm Exam(10/26)
    \item[9주차(10/28, 11/2)] Theta Grids and Functional Categories (Chapter 9)
    \item[10주차(11/4, 11/9)] Head-to-Head Movement (Chapter 10)
    \item[11주차(11/11, 11/16)] DP Movement (Chapter 11)
    \item[12주차(11/18, 11/23)] Wh-movement and Locality Constraints (Chapter 12)
    \item[13주차(11/25, 11/30)] A Unified Theory of Movement (Chapter 13)
    \item[14주차(12/2, 12/7)] Ditransitives (Chapter 14)
    \item[15주차(12/9)] Raising, Control, and Empty Categories (Chapter 15)
    \item[15주차(12/14)] Finalterm Exam 
\end{description}

\section{4. 평가방법}
\medskip
\begin{description}
    \item[\textbf{성적부여방식}] 상대평가(A--F)\footnote{S/U 선택가능}
    \item[\textbf{출석(10\%)}] 수업일수의 1/3을 초과하여 결석하면 성적은 ``F'' 또는 ``U''가 됨\\(담당교수가 불가피한 결석으로 인정하는 경우는 예외로 할 수 있음)
    \item[\textbf{과제(30\%)}] 퀴즈/과제 3회 
    \item[\textbf{중간(30\%)}] 
    \item[\textbf{기말(30\%)}] 
    \descriptionrule
    \item[\textbf{합계}] 100\%
\end{description}

\newpage
\section{5. 생성형 AI 도구 활용 방침}
\medskip
본 강의에서는 생성형 AI를 학습 보조 도구로 활용할 수 있습니다. 과제 수행에 필요한 아이디어 및 주제 탐색, 자료분석 등의 목적으로 생성형 AI를 사용할 수 있습니다. 다만, 생성형 AI를 과제 수행에 사용한 경우, 사용 도구명, 사용 과정 및 생성 내용 등을 과제에 명시해야 합니다.

\section{6. 정원 외 신청 수용 여부}
\medskip
최대 0명

\section{7. 장애학생 지원사항}
\medskip
\subsection{강의 수강 관련}
\medskip
\begin{itemize}
    \item 시각장애: 교재 제작(디지털교재, 점자교재, 확대교재 등), 대필도우미 허용
    \item 지체장애: 교재 제작(디지털교재), 대필도우미 및 수업보조 도우미 허용
    \item 청각장애: 대필 및 문자통역 도우미 활동 허용, 강의 녹취 허용
    \item 건강장애: 질병 등으로 인한 결석에 대한 출석 인정, 대필도우미 허용
    \item 학습장애: 대필도우미 허용
    \item 지적장애/자폐성장애: 대필도우미 및 수업 멘토 허용
\end{itemize}
\subsection{과제 및 평가 관련}
\medskip
\begin{itemize}
    \item 시각장애/지체장애/청각장애/건강장애/학습장애: 과제 제출기한 연장, 과제 제출 및 응답 방식의 조정, 평가 시간 연장, 평가 문항 제시 및 응답 방식의 조정, 별도 고사실 제공
    \item 지적장애/자폐성장애: 개별화 과제 제출 및 대체 평가 실시
\end{itemize}
\subsection{비고}
\medskip
본 강의를 수강하는 장애학생들에게는 이상의 지원 서비스 이외에도 장애학생 개개인의 특성과 요구에 따라, 지도교수 및 장애학생지원센터와의 상담을 통하여 적절한 수준의 지원 서비스를 제공합니다. 장애학생에 대한 지원서비스와 관련하여 문의사항이 있는 학생들은 담당교수 강초롱 혹은 장애학생지원센터(02-880-8787)로 문의바랍니다.

\end{document}
